%-------------------------------------------------------------------------------
% CAPITOLO 24
%-------------------------------------------------------------------------------
\chapter{Lezione 24}
\label{capitolo_24}
\textit{27/05/2026}

\section{Sequenze Proteiche e Modelli di Folding}
Non tutte le sequenze amminoacidiche producono la medesima struttura tridimensionale di una proteina. Ciononostante, si osserva che determinate sequenze, o famiglie di sequenze simili, ricorrono in natura. Questo fenomeno solleva interrogativi sulle proprietà che governano la selezione e la stabilità delle sequenze proteiche funzionali.



Consideriamo una famiglia di proteine:
\begin{itemize}
    \item ciascuna proteina è rappresentata da una sequenza di amminoacidi $\{S_i\}$
    \item una proteina tipica può avere una lunghezza $L=200$ amminoacidi $\longrightarrow i=1, \dots, L$
    \item ci sono 20 possibili amminoacidi $\longrightarrow S_i=1, \dots, 20$
\end{itemize}

Analizziamo l'entropia di queste sequenze per capire la statistica di occorrenza degli amminoacidi.
Il numero totale di sequenze possibili per una proteina di lunghezza $L=200$ è $N_{tot} = 20^{200}$.
L'entropia per una distribuzione di probabilità $P(S)$ su tutte le possibili sequenze è data da:

\begin{equation} \label{eq:ch24_1}
S = -\sum_s P(s) \ln[P(s)]
\end{equation}

dove la somma è su tutte le $20^L$ sequenze.

\subsection*{Assunzione più semplice: indipendenza dei siti}
Assumiamo che la scelta dell'amminoacido in una posizione $i$ sia indipendente dalla scelta degli amminoacidi nelle altre posizioni. La probabilità di una sequenza $S = (S_1, S_2, \dots, S_L)$ è allora il prodotto delle probabilità dei singoli amminoacidi in ciascuna posizione:

\begin{equation} \label{eq:ch24_2}
P(s) = \prod_{i=1}^{L} P^{(i)}(S_i)
\end{equation}

dove $P^{(i)}(S_i)$ è la probabilità che l'amminoacido $S_i$ si trovi in posizione $i$. Questa assunzione potrebbe essere errata.
Sotto questa ipotesi di indipendenza, l'entropia totale diventa additiva:

\begin{align} 
S &= -\sum_{\{S_j\}} \left( \prod_{k=1}^{L} P^{(k)}(S_k) \right) \ln\left[ \prod_{m=1}^{L} P^{(m)}(S_m) \right] \nonumber \\
&= -\sum_{S_1, \dots, S_L} \left( \prod_{k=1}^{L} P^{(k)}(S_k) \right) \sum_{m=1}^{L} \ln[P^{(m)}(S_m)] \nonumber \\
&= -\sum_{m=1}^{200} \Biggl( \sum_{S_m} P^{(m)}(S_m) \ln[P^{(m)}(S_m)] \Biggr) \nonumber \\
&= \sum_{i=1}^{L} S^{(i)} \label{eq:ch24_3}
\end{align}

dove $S^{(i)} = -\sum_{S_i} P^{(i)}(S_i) \ln[P^{(i)}(S_i)]$ è l'entropia associata al sito $i$.
Se non avessimo alcuna informazione, potremmo assumere una probabilità uniforme per ciascun tipo di amminoacido in ogni sito:

\begin{equation} \label{eq:ch24_4}
P^{(i)}(S_i) = \frac{1}{20} \quad \forall i, \forall S_i
\end{equation}

In questo caso, l'entropia per sito sarebbe:

\begin{equation} \label{eq:ch24_5}
S^{(i)} = -\sum_{S_i=1}^{20} \frac{1}{20} \ln\left(\frac{1}{20}\right) = \ln(20)
\end{equation}

E l'entropia totale per una proteina di lunghezza $L=200$ sarebbe:

\begin{equation} \label{eq:ch24_6}
S_{indip,unif} = 200 \ln(20)
\end{equation}

Tuttavia, analizzando un database di $M$ sequenze proteiche allineate, possiamo stimare le frequenze $P^{(i)}(S_i)$ per ogni sito $i$:

\begin{equation} \label{eq:ch24_7}
P^{(i)}(S_i=k) = \frac{n^{(i)}(k)}{M}
\end{equation}

dove $n^{(i)}(k)$ è il numero di volte che l'amminoacido di tipo $k$ appare al sito $i$ nelle $M$ sequenze.
I dati sperimentali ci dicono che non c'è omogeneità: le distribuzioni $P^{(i)}(S_i)$ variano da sito a sito. Quindi, l'assunzione di probabilità uniforme $1/20$ per ogni sito è errata.


\section{Assunzione oltre l'Indipendenza: Correlazioni tra Siti}
Se l'assunzione $P(s) = \prod_i P^{(i)}(S_i)$ non è corretta, dobbiamo considerare le dipendenze tra siti. Un primo passo è introdurre statistiche a due punti (correlazioni tra coppie di siti), come $P(S_i, S_j)$, la probabilità congiunta di osservare l'amminoacido $S_i$ al sito $i$ e l'amminoacido $S_j$ al sito $j$. Questa distribuzione vive in uno spazio di $20 \times 20 = 400$ stati per ogni coppia di siti.



Dai dati, possiamo stimare le correlazioni sperimentali:

\begin{equation} \label{eq:ch24_8}
C_{ij}^{exp} = \langle S_i S_j \rangle_{exp} - \langle S_i \rangle_{exp} \langle S_j \rangle_{exp} 
\end{equation}

L'idea è di costruire un modello probabilistico $P_{mod}(s)$ che riproduca non solo le frequenze dei singoli siti $P^{(i)}(S_i)$ ma anche le correlazioni osservate tra coppie di siti $P(S_i,S_j)$.
Si può usare un algoritmo Monte Carlo per generare un insieme di sequenze artificiali i cui profili di correlazione $C_{ij}^{mod}$ minimizzino la discrepanza con quelli sperimentali $C_{ij}^{exp}$:

\begin{equation} \label{eq:ch24_9}
\chi^2 = \sum_{i,j} |C_{ij}^{mod} - C_{ij}^{exp}|^2
\end{equation}

Studi mostrano che i modelli che includono le interazioni a due punti (correlazioni) sono significativamente migliori nel predire quali sequenze si ripiegano correttamente rispetto ai modelli basati sulla sola statistica dei singoli siti. Questo suggerisce che le interazioni (o co-evoluzioni) tra coppie di amminoacidi lungo la catena sono importanti per il folding. Queste coppie non sono necessariamente adiacenti nella sequenza primaria, ma possono essere in contatto nella struttura tridimensionale ripiegata.


\section{Modello di Massima Entropia (MEM)}
Un approccio formale per costruire modelli probabilistici che siano consistenti con un insieme di osservabili sperimentali è il Modello di Massima Entropia (MEM). Si cerca la distribuzione di probabilità $P(s)$ che massimizza l'entropia $S = -\sum_s P(s) \ln P(s)$ sotto il vincolo che le medie di certe osservabili $O^\mu(s)$ calcolate con $P(s)$ eguaglino i valori sperimentali $\langle O^\mu \rangle_{exp}$:

\begin{equation} \label{eq:ch24_10}
\sum_s P(s) O^\mu(s) = \langle O^\mu \rangle_{exp} \quad \forall \mu
\end{equation}

e la condizione di normalizzazione $\sum_s P(s) = 1$.
La soluzione a questo problema variazionale è una distribuzione della forma:

\begin{equation} \label{eq:ch24_11}
P_{MEM}(s) = \frac{1}{Z} \exp\left(+\sum_\mu \lambda_\mu O^\mu(s)\right)
\end{equation}

dove i $\lambda_\mu$ sono moltiplicatori di Lagrange determinati imponendo i vincoli; la funzione di partizione è data da:

\begin{equation} \label{eq:ch24_12}
Z = \sum_s \exp\left(+\sum_\mu \lambda_\mu O^\mu(s)\right)  
\end{equation}

La massimizzazione dell'entropia vincolata può essere formulata massimizzando il funzionale:

\begin{equation} \label{eq:ch24_13}
\tilde{S}[P, \{\lambda_\mu\}] = -\sum_s P(s) \ln P(s) + \sum_\mu \lambda_\mu \left( \sum_s P(s) O^\mu(s) - \langle O^\mu \rangle_{exp} \right)
\end{equation}

Sostituendo $P_{MEM}(s)$ nel funzionale dell'entropia, si ottiene:

\begin{align} 
\tilde{S}[P, \{\lambda_\mu\}] &= -\sum_s P(s) \Biggl[ -\ln(Z) + \left(\sum_\mu \lambda_\mu O^\mu(s)\right)\Biggr] \nonumber \\
&\quad + \sum_\mu \lambda_\mu \left( \sum_s P(s) O^\mu(s) - \langle O^\mu \rangle_{exp} \right) \label{eq:ch24_14}
\end{align}

Poichè $\sum_s P(s) \ln(Z) = 1 \cdot \ln(Z)$ a causa della normalizzazione, il risultato finale è:

\begin{equation} \label{eq:ch24_15}
\tilde{S}[\{\lambda_\mu\}] = \ln Z[\{\lambda_\mu\}] - \sum_\mu \lambda_\mu \langle O^\mu \rangle_{exp}
\end{equation}

Per determinare i valori ottimali dei $\lambda_\mu$, si pongono a zero le derivate parziali di $\tilde{S}$ rispetto a ciascun $\lambda_\mu$:

\begin{equation} \label{eq:ch24_16}
\frac{\partial\tilde{S}}{\partial\lambda_\mu} = 0
\end{equation}

Consideriamo ora la probabilità di osservare una specifica sequenza sperimentale $S_{exp}$ secondo un modello parametrizzato dai $\lambda_\mu$.

\begin{equation} \label{eq:ch24_17}
P(S_{exp}) = \frac{e^{\sum_\mu \lambda_\mu O^\mu(S_{exp})}}{Z}
\end{equation}

Questa è la probabilità di una singola realizzazione sperimentale $S_{exp}$. Se abbiamo $M$ realizzazioni sperimentali indipendenti $S_{exp}^{(\alpha)}$ (con $\alpha = 1, \dots, M$), indicando l'insieme di queste osservazioni come $\{S_{exp}\}$, la probabilità è:

\begin{equation} \label{eq:ch24_18}
P(\{S_{exp}^{(\alpha)}\}) = \prod_{\alpha=1}^{M} P(S_{exp}^{(\alpha)}) = \prod_{\alpha=1}^{M} \frac{e^{\sum_\mu \lambda_\mu O^\mu(S_{exp}^{(\alpha)})}}{Z}
\end{equation}

Prendendo il logaritmo naturale, si ottiene la Log-Likelihood:

\begin{equation} \label{eq:ch24_19}
\ln P(\{S_{exp}^{(\alpha)}\}) = \sum_{\alpha=1}^{M} \left( -\ln Z + \sum_\mu \lambda_\mu O^\mu(S_{exp}^{(\alpha)}) \right)
\end{equation}

Questa espressione può essere riscritta fattorizzando $M$:

\begin{equation} \label{eq:ch24_20}
\ln P(\{S_{exp}^{(\alpha)}\}) = M \left[ -\ln Z + \frac{1}{M} \sum_{\alpha=1}^{M} \sum_\mu \lambda_\mu O^\mu(S_{exp}^{(\alpha)}) \right]
\end{equation}

Utilizzando la definizione della media sperimentale per ogni osservabile, $\langle O^\mu \rangle_{exp} = \frac{1}{M} \sum_{\alpha=1}^{M} O^\mu(S_{exp}^{(\alpha)})$, la Log-Likelihood diventa:

\begin{equation} \label{eq:ch24_21}
\ln P(\{S_{exp}^{(\alpha)}\}) = M \left( -\ln Z + \sum_\mu \lambda_\mu \langle O^\mu \rangle_{exp} \right)
\end{equation}

Massimizzare questa Log-Likelihood rispetto ai $\lambda_\mu$ è equivalente a risolvere il sistema di equazioni $\langle O^\mu \rangle_{model} = \langle O^\mu \rangle_{exp}$.
